% Options for packages loaded elsewhere
\PassOptionsToPackage{unicode}{hyperref}
\PassOptionsToPackage{hyphens}{url}
%
\documentclass[
]{article}
\usepackage{amsmath,amssymb}
\usepackage{iftex}
\ifPDFTeX
  \usepackage[T1]{fontenc}
  \usepackage[utf8]{inputenc}
  \usepackage{textcomp} % provide euro and other symbols
\else % if luatex or xetex
  \usepackage{unicode-math} % this also loads fontspec
  \defaultfontfeatures{Scale=MatchLowercase}
  \defaultfontfeatures[\rmfamily]{Ligatures=TeX,Scale=1}
\fi
\usepackage{lmodern}
\ifPDFTeX\else
  % xetex/luatex font selection
\fi
% Use upquote if available, for straight quotes in verbatim environments
\IfFileExists{upquote.sty}{\usepackage{upquote}}{}
\IfFileExists{microtype.sty}{% use microtype if available
  \usepackage[]{microtype}
  \UseMicrotypeSet[protrusion]{basicmath} % disable protrusion for tt fonts
}{}
\makeatletter
\@ifundefined{KOMAClassName}{% if non-KOMA class
  \IfFileExists{parskip.sty}{%
    \usepackage{parskip}
  }{% else
    \setlength{\parindent}{0pt}
    \setlength{\parskip}{6pt plus 2pt minus 1pt}}
}{% if KOMA class
  \KOMAoptions{parskip=half}}
\makeatother
\usepackage{xcolor}
\usepackage[margin=1in]{geometry}
\usepackage{graphicx}
\makeatletter
\def\maxwidth{\ifdim\Gin@nat@width>\linewidth\linewidth\else\Gin@nat@width\fi}
\def\maxheight{\ifdim\Gin@nat@height>\textheight\textheight\else\Gin@nat@height\fi}
\makeatother
% Scale images if necessary, so that they will not overflow the page
% margins by default, and it is still possible to overwrite the defaults
% using explicit options in \includegraphics[width, height, ...]{}
\setkeys{Gin}{width=\maxwidth,height=\maxheight,keepaspectratio}
% Set default figure placement to htbp
\makeatletter
\def\fps@figure{htbp}
\makeatother
\setlength{\emergencystretch}{3em} % prevent overfull lines
\providecommand{\tightlist}{%
  \setlength{\itemsep}{0pt}\setlength{\parskip}{0pt}}
\setcounter{secnumdepth}{-\maxdimen} % remove section numbering
\usepackage{booktabs}
\usepackage{longtable}
\usepackage{array}
\usepackage{multirow}
\usepackage{wrapfig}
\usepackage{float}
\usepackage{colortbl}
\usepackage{pdflscape}
\usepackage{tabu}
\usepackage{threeparttable}
\usepackage{threeparttablex}
\usepackage[normalem]{ulem}
\usepackage{makecell}
\usepackage{xcolor}
\usepackage{caption}
\usepackage{anyfontsize}
\ifLuaTeX
  \usepackage{selnolig}  % disable illegal ligatures
\fi
\usepackage{bookmark}
\IfFileExists{xurl.sty}{\usepackage{xurl}}{} % add URL line breaks if available
\urlstyle{same}
\hypersetup{
  pdftitle={Hyper-Reflectivity and Neovascularization in LMRI DB},
  pdfauthor={Christian Anderson},
  hidelinks,
  pdfcreator={LaTeX via pandoc}}

\title{Hyper-Reflectivity and Neovascularization in LMRI DB}
\author{Christian Anderson}
\date{2025-02-25}

\begin{document}
\maketitle

\subsubsection{Q: What percent of MacTel patients have outer retina
hyper reflectivity (ORHR) and choroidal neovascularization
(CNV)?}\label{q-what-percent-of-mactel-patients-have-outer-retina-hyper-reflectivity-orhr-and-choroidal-neovascularization-cnv}

ORHR and CNV are not variables in the database, but HR\_TEMPORAL (Hyper
reflectivity on the temporal side) and NEOVASC (Neovascularization) are.

\textbf{36.6} \% of eyes develop HR\_TEMPORAL and \textbf{6.1\%} develop
NEOVASC while only \textbf{4.9\%} develop both.

\textbf{47.6} \% of patients develop HR\_TEMPORAL and \textbf{10.1\%}
develop NEOVASC while only \textbf{8.7\%} develop both.

These percentages are found in Table 2 and 4, respectively.

If you are interested in seeing a breakdown of phenotype presence on the
level of patients and eyes, please reference the tables below.

\paragraph{Tables}\label{tables}

\begin{table}[!t]
\caption*{
{\large Eyes that develop NEOVASC and/or HR\_TEMPORAL} \\ 
{\small 0: not Present; 1: Questionable; 2: Present}
} 
\fontsize{12.0pt}{14.4pt}\selectfont
\begin{tabular*}{450pt}{@{\extracolsep{\fill}}l|l|rrr}
\toprule
\multicolumn{2}{l}{} & \multicolumn{3}{c}{HR\_TEMPORAL} \\ 
\cmidrule(lr){3-5}
\multicolumn{2}{c}{} & 0 & 1 & 2 \\ 
\midrule\addlinespace[2.5pt]
\multirow{3}{*}{NEOVASC} & 0 & 4535 & 1039 & 1266 \\ 
 & 1 & 63 & 47 & 120 \\ 
 & 2 & 21 & 20 & 174 \\ 
\bottomrule
\end{tabular*}
\begin{minipage}{\linewidth}
X\textasciicircum{}2 test: p \textasciitilde{} 1.061e-138; N = 7285\\
Note: These are counts of unique eyes that ever develop phenotype, not counts of eye-visits\\
\end{minipage}
\end{table}

\textbf{Table 1}: Counts of eyes for every combination of HR\_TEMPORAL
and NEOVASC. A chi-square test was performed on these values to test
whether the presence of Neovascularization on Hyper-reflectivity and
vice-versa.

\begin{center}\rule{0.5\linewidth}{0.5pt}\end{center}

\begin{center}\rule{0.5\linewidth}{0.5pt}\end{center}

\begin{table}[!t]
\caption*{
{\large Ratio of eyes that develop NEOVASC and/or HR\_TEMPORAL} \\ 
{\small 0: not Present; 1: Questionable; 2: Present}
} 
\fontsize{12.0pt}{14.4pt}\selectfont
\begin{tabular*}{450pt}{@{\extracolsep{\fill}}l|l|rrr}
\toprule
\multicolumn{2}{l}{} & \multicolumn{3}{c}{HR\_TEMPORAL} \\ 
\cmidrule(lr){3-5}
\multicolumn{2}{c}{} & 0 & 1 & 2 \\ 
\midrule\addlinespace[2.5pt]
\multirow{3}{*}{NEOVASC} & 0 & 0.623 & 0.143 & 0.174 \\ 
 & 1 & 0.009 & 0.006 & 0.016 \\ 
 & 2 & 0.003 & 0.003 & 0.024 \\ 
\bottomrule
\end{tabular*}
\begin{minipage}{\linewidth}
Ratio develop both: 0.049000 (1,1,+1,2+2,2+2,1)\\
Ratio develop HR\_TEMPORAL only: 0.317000 (0,1+0,2)\\
Ratio develop NEOVASC only: 0.012000 (1,0+2,0)\\
\end{minipage}
\end{table}

\textbf{Table 2}: Every combination of HR\_TEMPORAL and NEOVASC as a
ratio of total eyes (`N' in Table 1). The percentages of eyes that
develop both and each individually were calculated by summing their
respective cells in the table.

\begin{center}\rule{0.5\linewidth}{0.5pt}\end{center}

\begin{center}\rule{0.5\linewidth}{0.5pt}\end{center}

\begin{table}[!t]
\caption*{
{\large Patients that develop NEOVASC and/or HR\_TEMPORAL} \\ 
{\small 0: not Present; 1: Questionable; 2: Present}
} 
\fontsize{12.0pt}{14.4pt}\selectfont
\begin{tabular*}{450pt}{@{\extracolsep{\fill}}l|l|rrr}
\toprule
\multicolumn{2}{l}{} & \multicolumn{3}{c}{HR\_TEMPORAL} \\ 
\cmidrule(lr){3-5}
\multicolumn{2}{c}{} & 0 & 1 & 2 \\ 
\midrule\addlinespace[2.5pt]
\multirow{3}{*}{NEOVASC} & 0 & 1869 & 614 & 810 \\ 
 & 1 & 36 & 40 & 106 \\ 
 & 2 & 14 & 14 & 158 \\ 
\bottomrule
\end{tabular*}
\begin{minipage}{\linewidth}
X\textasciicircum{}2 test: p \textasciitilde{} 9.332e-87; N = 3661\\
Note: These are counts of unique patients that ever develop phenotype, not counts of patient-visits\\
\end{minipage}
\end{table}

\textbf{Table 3}: Counts of patients for every combination of
HR\_TEMPORAL and NEOVASC. A chi-square test was performed on these
values to test whether the presence of Neovascularization on
Hyper-reflectivity and vice-versa.

\begin{center}\rule{0.5\linewidth}{0.5pt}\end{center}

\begin{center}\rule{0.5\linewidth}{0.5pt}\end{center}

\begin{table}[!t]
\caption*{
{\large Ratio of patients that develop NEOVASC and/or HR\_TEMPORAL} \\ 
{\small 0: not Present; 1: Questionable; 2: Present}
} 
\fontsize{12.0pt}{14.4pt}\selectfont
\begin{tabular*}{450pt}{@{\extracolsep{\fill}}l|l|rrr}
\toprule
\multicolumn{2}{l}{} & \multicolumn{3}{c}{HR\_TEMPORAL} \\ 
\cmidrule(lr){3-5}
\multicolumn{2}{c}{} & 0 & 1 & 2 \\ 
\midrule\addlinespace[2.5pt]
\multirow{3}{*}{NEOVASC} & 0 & 0.511 & 0.168 & 0.221 \\ 
 & 1 & 0.010 & 0.011 & 0.029 \\ 
 & 2 & 0.004 & 0.004 & 0.043 \\ 
\bottomrule
\end{tabular*}
\begin{minipage}{\linewidth}
Ratio develop both: 0.087000 (1,1,+1,2+2,2+2,1)\\
Ratio develop HR\_TEMPORAL only: 0.389000 (0,1+0,2)\\
Ratio develop NEOVASC only: 0.014000 (1,0+2,0)\\
\end{minipage}
\end{table}

\textbf{Table 4}: Every combination of HR\_TEMPORAL and NEOVASC as a
ratio of total patients (`N' in Table 3). The percentages of patients
that develop both and each individually were calculated by summing their
respective cells in the table.

\end{document}
